\noindent\textbf{Elección del procedimiento.} Con el convenio de circulación que produce sustentación positiva, Kutta--Joukowski da $L\prime=\rho V\Gamma$. Dividir por $q c$ cancela densidad y deja $c_l=2\Gamma/(Vc)$. Si se adopta la aproximación de perfil delgado $c_l\simeq2\pi\alpha$, el ángulo sale inicialmente en radianes.

\begin{enumerate}

\item La fuerza corresponde a unidad de envergadura.

\[\fitmath{\mathrm{Lp}=\rho\cdot V\cdot G}\]

\[\fitmath{\mathrm{Lp}=(1.225)\cdot (60)\cdot (40)\approx 2940\,\mathrm{N/m}}\]

\item Usar cuerda como longitud de referencia.

\[\fitmath{\mathrm{cl}=\frac{2\cdot G}{V\cdot c}}\]

\[\fitmath{\mathrm{cl}=\frac{2\cdot (40)}{(60)\cdot (1.5)}\approx 0.8888889\,}\]

\item Dividir por la pendiente por radián y después convertir a grados.

\[\fitmath{\alpha=\frac{\frac{\mathrm{cl}}{2\cdot \pi}\cdot 180}{\pi}}\]

\[\fitmath{\alpha=\frac{\frac{(0.8888888889)}{2\cdot \pi}\cdot 180}{\pi}\approx 8.105695\,{}^{\circ}}\]

\end{enumerate}
